\begin{proof}[Proof of \cref{obj:lem:localized-vc-process-bound}]
The policy-class, margin, zero-effect, well-formedness, and bounded-outcome conditions delimit the setting of the lemma; no quantitative use is made of them here.  The displayed empirical-process inequality itself is the fixed-radius localized envelope supplied by \cref{obj:ass:vc-localized-envelope}.  That assumption gives constants \(C>0\) and \(p\ge0\) such that, for every sample size \(m\ge1\), envelope \(B\ge0\), radius \(r\ge0\), and policy-compatible increment family \(g_\pi\) satisfying
\[
 |g_\pi(O)|\le B\quad(\pi\in\Pi)
 \qquad\text{and}\qquad
 \int g_\pi(O)^2\,dP(O)\le B^2P_X(D_\pi)\quad(\pi\in\Pi),
\]
the expected supremum of the centered empirical process over the regret-localized class \(\{\pi\in\Pi:R_P(\pi)\le r\}\) is bounded at the stated rate.  Its localized supremum is exactly the quantity named in the lemma: for \(O_1,\ldots,O_m\) drawn i.i.d.\ from \(P\),
\[
Z_m(r)
=
\sup_{\pi\in\Pi:\,R_P(\pi)\le r}\left|(P_m-P)g_\pi\right|
=
\sup_{\pi\in\Pi:\,R_P(\pi)\le r}
\left|m^{-1}\sum_{i=1}^m g_\pi(O_i)-\int g_\pi(O)\,dP(O)\right| ,
\]
so the assumption yields, with the same constants \(C\) and \(p\),
\[
\mathbb E_P Z_m(r)
\le
C B m^{-1/2}r^{\alpha/(2+2\alpha)}(\log m)^p
=
C B m^{-1/2}r^{\kappa}(\log m)^p ,
\]
the last equality being the definition \(\kappa=\alpha/(2+2\alpha)\) in the statement.  No separate symmetrization, chaining, or derivation from the other standing hypotheses is used in this lemma.
% lean: VCLocalizedEnvelope % lean: localized_vc_process_bound
\end{proof}
